\subsection{X/H 仿紧的情形}

若 X/H 是仿紧的，我们先来考察：向量空间 \(\mathcal{K}^{\chi}(X)\)（其中 \(\chi\) 变动）彼此同构，特别地都与 \(\mathcal{K}^{1}(X)\) 同构。

命题 7. --- 假设 X/H 仿紧。设 \(\chi\) 是 H 到 \(\mathbf{R}_{+}^{*}\) 的连续表示。

\begin{enumerate}
\def\labelenumi{\alph{enumi})}
\item
  在 \(\mathbf{X}\) 上存在连续函数 \(r\)，取值为 \(>0\)，满足 \(r(x\xi) = \chi (\xi)r(x)\)，对所有 \(x\in \mathbf{X}\) 和 \(\xi \in \mathbf{H}\) 成立。
\item
  映射 \(g \mapsto g / r\) 是从向量空间 \(\mathcal{H}^{\chi}(\mathbf{X})\) 到向量空间 \(\mathcal{H}^{1}(\mathbf{X})\) 的同构。
\end{enumerate}

应用 No.~1 的 Prop. 1，取 f 为一个 \(\geqslant0\) 且在任何轨道上不恒为零的函数（由附录 1 的引理 1 可做到）；则 \(r = f^{x}\) 满足 a) 的性质。b) 的断言显然。

命题 8. --- 假设 X/H 仿紧。在 X 上存在连续函数 \(h \geqslant 0\)，其支集与 X 的每个紧子集的饱和集的交为紧，并满足 \(h^{b}=1\)。对于这样的函数，X/H 上每个连续函数 g 都满足 \(g=\left(h\cdot(g\circ\pi)\right)^{b}\)。

应用 No.~1 的 Prop. 1，取 \(\chi = 1\)，并取 \(f\) 为一个 \(\geqslant 0\) 且在任何轨道上不恒为零的函数。我们有 \(f^1(x) > 0\)，对 X 的每个点 \(x\) 成立。置 \(h = f / f^1\)。于是 \(h^1 = f^1 / f^1 = 1\)，从而 \(h^b = 1\)。因此，若 \(g\) 是 X/H 上的连续函数，则 \(\left(h \cdot (g \circ \pi)\right)^b = h^b \cdot g = g\)。

注记。--- 1) 特别地，设 X 是一个局部紧空间，离散群 D 在其上连续且 逆紧地从右作用；假设 X/D 仿紧。则 X 上存在连续函数 \(h \geqslant 0\)，其支集与 X 的每个紧子集的饱和集的交为紧，并且 \(\sum_{d \in D} h(x d) = 1\) 对每个 \(x \in X\) 成立（除有限项外，该和的所有项均为零）。

\begin{enumerate}
\def\labelenumi{\arabic{enumi})}
\setcounter{enumi}{1}
\tightlist
\item
  保持 Prop. 8 的假设和记号。映射 \(g \mapsto h \cdot (g \circ \pi)\) 是从 \(\mathcal{K}(\mathrm{X / H})\) 到 \(\mathcal{K}(\mathrm{X})\) 的连续映射，并且是映射 \(f \mapsto f^{\flat}\) 的右逆。因此，\(\mathcal{K}(\mathrm{X / H})\) 的每个有界（分别为紧）子集都是 \(\mathcal{K}(\mathrm{X})\) 的某个有界（分别为紧）子集的像。由此立即得到，当这些空间赋予有界（分别为紧）收敛的拓扑时，映射 \(\lambda \mapsto \lambda^{\sharp}\) 是从 \(\mathcal{M}(\mathrm{X / H})\) 到 \(\mathcal{M}(\mathrm{X})\) 的一个闭线性子空间上的同构。
\end{enumerate}

命题 9. --- 保持 Prop. 8 的假设和记号。设 \(\lambda\) 是 X/H 上的正测度。

\begin{enumerate}
\def\labelenumi{\alph{enumi})}
\item
  偶 \((\pi, h)\) 对 \(\lambda^{\sharp}\) 而言是适合的，并且 \(\int_{\mathbf{X}} h(x) \varepsilon_{\pi(x)} d\lambda^{\sharp}(x) = \lambda\)。
\item
  映射 \(\pi\) 关于测度 \(h\cdot \lambda^{\sharp}\) 是 逆紧的，并且 \(\pi (h\cdot \lambda^{\sharp}) = \lambda\)。
\item
  设 \(\mathbf{k}\) 是 \(\mathrm{X / H}\) 上取值于 Banach 空间或 \(\overline{\mathbf{R}}\) 的函数。\(\mathbf{k}\) 关于 \(\lambda\) 可测（分别为局部可积、本质可积、可积）的充要条件是 \(h\cdot (\mathbf{k}\circ \pi)\) 关于 \(\lambda^{\sharp}\) 具有相应性质；并且，若 \(\mathbf{k}\) 关于 \(\lambda\) 本质可积，则
\end{enumerate}

\[
\int_ {\mathrm{X} / \mathrm{H}} \mathbf {k} d \lambda = \int_ {\mathrm{X}} h \cdot (\mathbf {k} \circ \pi) d \lambda^ {\sharp}.\tag{14}
\]

设 \(f \in \mathcal{K}(\mathrm{X}/\mathrm{H})\)。则 \(h \cdot (f \circ \pi) \in \mathcal{K}(\mathrm{X})\)，并且

\[
\int_ {\mathrm{X}} h (x) f (\pi (x)) d \lambda^ {\sharp} (x) = \int_ {\mathrm{X} / \mathrm{H}} f (\dot {x}) d \lambda (\dot {x}) \int_ {\mathrm{H}} h (x \xi) d \beta (\xi) = \int_ {\mathrm{X} / \mathrm{H}} f (\dot {x}) d \lambda (\dot {x}),
\]

由此得 a)。b) 的断言类似可证。关于 c) 的可测性、本质可积性和公式 (14) 的断言，随后可应用 Ch. V（§4, Prop. 3、§5, Th. 1、§4, Th. 2）的结果得到。若 k 是 λ-可积的，则 \(h \cdot (\mathbf{k} \circ \pi)\) 是 \(\lambda^{\sharp}\)-可积的（Ch. V, §3, No.~3, Th. 1）。若 \(h \cdot (\mathbf{k} \circ \pi)\) 是 \(\lambda^{\sharp}\)-可积的，则 Prop. 5 证明 \((h \cdot (\mathbf{k} \circ \pi))^b = h^b \cdot \mathbf{k} = \mathbf{k}\) 是 λ-可积的。若 k 是局部 λ-可积的，则 \(h \cdot (\mathbf{k} \circ \pi)\) 是局部 \(\lambda^{\sharp}\)-可积的（Prop. 6）。最后，假设 \(h \cdot (\mathbf{k} \circ \pi)\) 局部 \(\lambda^{\sharp}\)-可积；对于每个 \(f \in \mathcal{H}(\mathrm{X}/\mathrm{H})\)，\(h \cdot (\mathbf{k} \circ \pi) \cdot (f \circ \pi)\) 有紧支集，并且

\[
\left| h \cdot (\mathbf {k} \circ \pi) \cdot (f \circ \pi) \right| \leqslant \mathrm{M} \left| h \cdot (\mathbf {k} \circ \pi) \right|,
\]

其中 \(M = \sup |f|\)；因此 \(h \cdot ((\mathbf{k}f) \circ \pi)\) 是 \(\lambda^{\sharp}\)-可积的，从而根据已经证明的结果，kf 是 \(\lambda\)-可积的；这就证明了 k 确实是局部 \(\lambda\)-可积的。

推论。--- Prop. 5 所定义的映射 \(f \mapsto f^{b}\)，从 \(\mathrm{L}^{1}(\mathrm{X}, \lambda^{\sharp})\) 到 \(\mathrm{L}^{1}(\mathrm{X}/\mathrm{H}, \lambda)\)，是连续线性满射。

先假设 X/H 仿紧，并设 h 是 X 上满足 Prop. 8 条件的函数。若 k 是 X/H 上的 \(\lambda\)-可积数值函数，则 \(h \cdot (k \circ \pi)\) 是 \(\lambda^{\sharp}\)-可积的，且 \((h \cdot (k \circ \pi))^{\flat} = k\)（Prop. 9）。

一般情形下，令 \(u \in \mathrm{L}^1(\mathrm{X/H},\lambda)\)。存在一个 \(f \in \mathcal{L}^1(\mathrm{X/H},\lambda)\)，其等价类为 \(u\)，并且在紧集序列 \(\mathbf{K}_n\) 的可数并之外为零。递归地定义相对紧开集序列 \(\mathrm{U}_n\)，它们位于 \(\mathrm{X/H}\) 中，使得 \(\mathrm{U}_{n+1} \supset \mathrm{K}_n \cup \overline{\mathrm{U}}_n\)，并令 \(V\) 为 \(\mathrm{U}_n\) 的并。则 \(V\) 是 \(\mathrm{X/H}\) 的开子集，是紧子集 \(\overline{\mathrm{U}}_n\) 的可数并，因而是仿紧的（GT, I, §9, No.~10, Th. 5）。置 \(Y = \overline{\pi}^1(V)\)，并令 \(\lambda_V\)（分别为 \(\lambda_Y^\sharp\)）表示 \(\lambda\)（分别为 \(\lambda^\sharp\)）在 \(V\)（分别在 \(Y\)）上诱导的测度。显然，\(Y/H\) 可与 \(V\) 识别（GT, I, §3, Prop. 10），且 \(\lambda_Y^\sharp\) 可与 \((\lambda_V)^{\sharp}\) 识别。此外，\(f\) 在 \(V\) 外为零，并属于 \(\mathcal{L}^1(V,\lambda_V)\)。因此，存在 \(g \in \mathcal{L}^1(Y,\lambda_Y^\sharp)\)，使得 \(g^b = f\) 在 \(V\) 中几乎处处成立。将 \(g\) 在 \(X - Y\) 上延拓为 0，得到函数 \(g_1 \in \mathcal{L}^1(X,\lambda^\sharp)\)，显然，\(g_1^b\) 在 \(L^1(X/H,\lambda)\) 中的等价类正是 \(u\)。

注记 3. --- 假设 X/H 仿紧，并保留命题 9 的记号。此时，由 (14)，映射 \(k \mapsto h \cdot (k \circ \pi)\) 从 \(L^1(X/H, \lambda)\) 到 \(L^1(X, \lambda^\sharp)\) 是等距的，并且是映射 \(f \mapsto f^\flat\) 从 \(L^1(X, \lambda^\sharp)\) 到 \(L^1(X/H, \lambda)\) 的右逆。

